\documentclass{article}

\usepackage{amsmath}
\usepackage{tabularx}
\usepackage{booktabs}
\usepackage{longtable}
\usepackage{hyperref}
\hypersetup{
    colorlinks=true,
    linkcolor=blue,
    urlcolor=blue,
    citecolor=blue
}

\title{Reflection and Traceability Report on \progname}

\author{\authname}

\date{}

\input{../Comments}
\input{../Common}

\begin{document}

\maketitle

\tableofcontents

\newpage

\section{Changes in Response to Feedback}

This section documents every item of feedback received from TAs, the course
instructor, peer reviewers, the project supervisor, and usability testers
over the course of the project.  For each item the source, a short label,
and the change made (or the reason no change was made) are recorded.

The team's external supervisor is Luke, the MES Infra-Tech representative,
who attended regular bi-weekly stakeholder check-ins throughout the project.
The course instructor (Prof.\ Smith) observed the Rev~0 demonstration and
provided written feedback reproduced in part below.

\subsection{SRS and Hazard Analysis}

\begin{longtable}{p{3.2cm} p{3.8cm} p{7.4cm}}
\toprule
\textbf{Source} & \textbf{Feedback Label} & \textbf{Response / Change Made} \\
\midrule
\endhead

TA Meeting \#2 (Oct~2, 2025)
  & Requirements traceability matrix missing
  & Added traceability matrices to the SRS linking every functional and
    non-functional requirement to the corresponding VnV test case and, where
    applicable, to hazard analysis safety requirements.
    See PR~\#35 (\textit{Initial traceability matrices}).
\\[4pt]

TA Meeting \#2 (Oct~2, 2025)
  & Formal mathematical notation (SpecMath) required
  & Added a finite-state machine (FSM) representation of the user
    reimbursement workflow to the SRS to satisfy the formal-notation
    requirement.
\\[4pt]

TA Meeting \#2 (Oct~2, 2025)
  & Use case diagram recommended for system/technical section
  & Added a use case diagram to the SRS system-description section
    illustrating interactions between the three user roles (club member,
    club executive, VP Finance/MES admin) and the system.
\\[4pt]

TA Meeting \#2 (Oct~2, 2025)
  & FMEA table not traced back to SRS requirements
  & Added a ``Related SRS Requirements'' column to the FMEA table in
    \texttt{HazardAnalysis.tex} so every hazard entry links to the safety
    or functional requirements that mitigate it.
\\[4pt]

TA Meeting \#2 (Oct~2, 2025)
  & No applicable privacy/security standards identified
  & Identified relevant data-handling and privacy standards for student
    financial data and added them to the SRS and Hazard Analysis.
\\[4pt]

Peer review (GitHub Issues)
  & Hazard categorisation ambiguity as technical and safety overlap
  & Clarified hazard categorisation criteria in
    \texttt{HazardAnalysis.tex}; hazards that produce both technical and
    safety consequences are now cross-referenced with a note, and the FMEA
    table was reviewed for consistent categorisation throughout.
\\[4pt]

Peer review (GitHub Issues)
  & Mitigation strategies too vague since no implementation specifics
  & Expanded the mitigation-strategy column of the FMEA table with
    concrete countermeasures for each system component (Authentication,
    Request Submission, Budget Tracking, Notifications, Admin Tools,
    Data Storage).
\\[4pt]

Peer review (GitHub Issues)
  & Safety-requirement implementation timeline unclear (SR1--SR10
    ``first'' vs.\ ``future'' lacks justification)
  & Added rationale for implementation prioritisation in the Safety
    Requirements section, explaining why certain requirements (SR2,
    SR4--SR8) were implemented in the initial release and others (SR1,
    SR3, SR9--SR10) were deferred to a later phase.
\\[4pt]

Peer review (GitHub Issues)
  & Missing/unverifiable citation (IppolitoWallace1995)
  & Verified that the reference file path
    (\texttt{../../refs/References}) is correct and that the full
    citation is present in the bibliography.
\\

\bottomrule
\end{longtable}

\subsection{Design and Design Documentation}

\begin{longtable}{p{3.2cm} p{3.8cm} p{7.4cm}}
\toprule
\textbf{Source} & \textbf{Feedback Label} & \textbf{Response / Change Made} \\
\midrule
\endhead

Peer review GitHub Issue \#81
  & Inconsistency in MIS Module M4
  & Resolved the inconsistency in Module~M4 (Receipt Processing Module)
    in \texttt{MIS.tex}.
    Commit: \texttt{3fc9799}
    (\textit{Address issue 81 - Inconsistency in MIS M4}).
\\[4pt]

Peer review GitHub Issue \#84
  & Consistency issues across MIS
  & Fixed consistency issues (naming, notation, table style) throughout
    \texttt{MIS.tex}.
    Commit: \texttt{f74a647}
    (\textit{Fix issue 84 - Consistency}).
\\[4pt]

Peer review GitHub Issue \#85
  & Receipt Processing Module (M4) misclassified
  & Reclassified M4 to the correct module category in
    \texttt{MIS.tex}.
    Commit: \texttt{8d2581d}
    (\textit{Issue 85: The Receipt Processing Module (M4) appears to be
    misclassified}).
\\[4pt]

Peer review GitHub Issue \#87
  & Symbols, Abbreviations and Acronyms section incomplete
  & Expanded the Symbols/Abbreviations/Acronyms table in
    \texttt{MIS.tex} with all missing terms.
    Commit: \texttt{5c84f06}
    (\textit{Issue 87: Incomplete Symbols, Abbreviations and Acronyms
    Section}).
\\[4pt]

Peer review GitHub Issue \#88
  & Missing type definitions for Parameters, Constants, and State
    Variables
  & Added explicit type definitions for all module parameters, constants,
    and state variables throughout \texttt{MIS.tex}.
    Commit: \texttt{c2ad946}
    (\textit{Issue 88: Missing Type Definitions}).
\\[4pt]

Design doc marking (TA/grader)
  & MIS table formatting and notation unclear; template text still
    present
  & Added table grid/wrapping fixes (\texttt{makecell}, custom column
    types), introduced formal map notation
    ($\text{map}\; A \rightarrow B$), clarified the Sequence/String/
    Tuple/Map type definitions, and removed all \texttt{\textbackslash
    wss\{\}} placeholder text.
    Commit: \texttt{8b695e7} (PR~\#94
    \textit{Add feedback from design doc marking}).
\\[4pt]

Design doc marking (TA/grader)
  & Anticipated Change (AC) description for external API integrations
    too vague
  & Updated the AC entry to read ``External API integrations with
    required services and actions (e.g., payment processing, email
    services)'', making explicit what services and actions are subject to
    change.
    Commit: \texttt{1ec121e} (PR~\#96 \textit{Feedback Implementation}).
\\[4pt]

Rev~0 demo Prof.\ Smith
  & Information hiding and likely changes not clearly explained during
    Q\&A
  & The MG already contained AC/UC sections; following the Rev~0
    feedback, the justification for each anticipated change was expanded
    to explain \textit{why} that decision was hidden and what a
    substitution would require.
\\[4pt]

Rev~0 demo Prof.\ Smith
  & Demo should be shown before discussing design
  & For the final presentation the ordering was adjusted so that the live
    system demo was shown first to ground the design discussion in
    observable behaviour.
\\[4pt]

Rev~0 demo Prof.\ Smith
  & Show more of the actual application rather than generic graphics
  & Final presentation slides were updated to include screenshots and
    screen recordings from the live system rather than placeholder
    diagrams.
\\

\bottomrule
\end{longtable}

\subsection{VnV Plan and Report}

\begin{longtable}{p{3.2cm} p{3.8cm} p{7.4cm}}
\toprule
\textbf{Source} & \textbf{Feedback Label} & \textbf{Response / Change Made} \\
\midrule
\endhead

TA Meeting \#3 (Oct~23, 2025)
  & VnV test cases not explicitly linked to SRS FR/NFR numbers
  & System test cases were rewritten to reference SRS requirement
    identifiers directly
    (\texttt{FRQ-1-SUBMIT} through \texttt{FRQ-6}, \texttt{NFL-SAL-1},
    \texttt{NFL-USAB-1}, \texttt{NFL-SEC-1}, \texttt{NFL-MNT-1}).
    PRs \#48, \#50--\#52.
\\[4pt]

TA Meeting \#3 (Oct~23, 2025)
  & Distinction between automated and manual tests not stated
  & Each test in the VnV Plan now has an explicit method tag
    (automated: UI automation, API check, database query; or
    manual: role-based UI walkthrough).
    PRs \#113--\#114.
\\[4pt]

TA Meeting \#3 (Oct~23, 2025)
  & Software Validation section (3.7) incomplete
  & Completed the Software Validation section to document how the team
    verified that the system meets the client's actual needs (MES
    stakeholder sign-off), distinct from technical verification.
    PR~\#115.
\\[4pt]

TA Meeting \#3 (Oct~23, 2025)
  & VnV template version outdated
  & Updated to the current VnV template (v2.4).
\\[4pt]

Rev~0 demo Prof.\ Smith / TA
  & Usability testing lacked a structured plan: no predefined
    investigation questions, interaction script, or measurement
    instrument
  & Rewrote the usability testing plan to include:
    (a) specific questions and concerns to investigate,
    (b) a structured five-task list for participants,
    (c) a think-aloud protocol,
    (d) defined measurement outcomes.
    The revised plan was executed with the MES Infratech Team on
    April~5, 2026 (see
    \texttt{docs/Extras/UsabilityTesting/UsabilityTesting.tex}).
\\[4pt]

Rev~0 demo Prof.\ Smith / TA
  & CI/CD plan amounts to automatic PDF generation only
  & Expanded the CI/CD section of the VnV Plan to cover ESLint linting,
    TypeScript type checking, and deployment to the MES staging branch,
    in addition to the LaTeX PDF generation pipeline.
\\[4pt]

Rev~0 demo TA
  & Instructor annotation comments (\texttt{\textbackslash wss\{\}})
    left in submitted PDFs
  & Removed all \texttt{\textbackslash wss\{\}} annotations from all
    deliverable source files.
    Commit: PR~\#123 (\textit{remove pink wss comments}).
\\[4pt]

Rev~0 demo TA
  & Team productivity report inaccurate
    (wrong lecture and TA meeting counts)
  & Corrected all attendance figures in the team contribution documents
    and committed accurate data.  All future contribution reports are
    verified against calendar records before submission.
\\[4pt]

PRs \#111--\#117
  & VnV Report sections incomplete at Rev~0
  & Completed all outstanding sections of \texttt{VnVReport.tex}:
    functional requirements evaluation (Sections 3--4, PR~\#111),
    unit testing section with manual-first strategy and tooling
    aligned to the MES repo (PR~\#114), appendix reflections
    (PR~\#115), and unit-testing alignment cleanup (PR~\#117).
\\

\bottomrule
\end{longtable}

% ----------------------------------------------------------------
\section{Extras}

\subsection{Usability Testing}

Usability testing was proposed as an extra deliverable at the POC demo
(November~20, 2025) and approved by the course instructor.  A formal
session was conducted on April~5, 2026 with the MES Infratech Team, who
represented the primary administrative users of the system and participated
with no prior documentation or training.

The session used an open-ended, task-based format.  Five representative tasks were evaluated: (1)~submitting a
receipt, (2)~approving/rejecting expense requests as a club executive,
(3)~viewing expenses filtered by year, (4)~accepting/rejecting requests as
VP Finance, and (5)~general workflow and navigation.

Key findings included moderate discoverability friction for receipt
submission (buried under Administration $\rightarrow$ Finances), an
expectation of drag-and-drop file upload that the system did not support, a
state-management bug causing auth-locked pages to appear empty on first
visit, a non-intuitive hover-only ``Details'' button on the VP Finance
dashboard, and a missing receipt-visibility feature for VP Finance
reviewers during approval.

Overall participant reception was positive: the system was described as a
significant improvement over the existing manual process and as ``the
budding of an ecosystem.''  Full findings and prioritised recommendations
are documented in
\texttt{docs/Extras/UsabilityTesting/UsabilityTesting.tex}.

\subsection{User Instructional Video}

A user instructional video was created as an approved extra, replacing a
traditional user manual as suggested by Prof.\ Smith at the POC demo.  The
video teaches MES club treasurers and administrators how to submit receipts,
review and approve reimbursement requests, and use the VP Finance dashboard.
It was added to the repository on April~21, 2026 and is accessible at:

\begin{center}
\url{https://youtu.be/toWvJ0CXTSU}
\end{center}

% ----------------------------------------------------------------
\section{Design Iteration (LO11 (PrototypeIterate))}

The design of the MES Club Payment Tracking System evolved through four
recognisable phases, each shaped by a combination of technical findings,
client requirements, and user feedback.

\paragraph{Initial design.}
The original concept was a form-based receipt submission interface that
would centralise the manual Google~Forms and spreadsheet workflow used by
MES clubs.  At this stage there was no automated data extraction; the
assumption was that users would manually enter all financial fields.  The
key architectural decision from the outset was integration into the existing
MES monorepo (\texttt{MES-Website-App-Router}) using the slices
architecture, rather than building a standalone application.

\paragraph{After the Proof-of-Concept (November 2025).}
The most significant technical risk was automated extraction of amounts,
line items, and tax from receipt photographs.  The POC resolved this by
demonstrating a viable pipeline using the Google Gemini LLM API, which
extracted structured data accurately across diverse receipt formats.  With
this risk eliminated, the design shifted focus to the remaining high-risk
item: user authentication and role-based access control (RBAC).

\paragraph{After Revision~0 (January--March 2026).}
RBAC was implemented with three permission tiers: club member (submit
only), club executive (submit + review club submissions), and VP
Finance/MES admin (system-wide review and approval).  Prof.\ Smith's Rev~0
feedback noted that the dashboard needed more realistic data for the
demonstration to be convincing, and that user testing was essential for the
project to demonstrate adoption readiness.  These observations directly
shaped the decision to conduct structured usability testing with actual MES
administrators rather than relying on internal testing alone.

\paragraph{After usability testing (April 2026).}
The MES Infratech Team session identified four high-priority UI issues that
were subsequently addressed:
\begin{itemize}
  \item \textbf{Receipt submission discoverability:} navigation depth
    reduced by adding a shortcut from the home page.
  \item \textbf{Auth-locked empty page:} a state-management bug was fixed
    so that protected pages hydrate correctly on first visit.
  \item \textbf{Details button affordance:} the hover-only interaction was
    replaced with a clearly styled clickable element.
  \item \textbf{Receipt visibility for VP Finance:} the review interface
    was updated to display the uploaded receipt alongside request details.
\end{itemize}

% ----------------------------------------------------------------
\section{Design Decisions (LO12)}

\subsection{Limitations}

\begin{itemize}
  \item \textbf{Google Gemini free-tier rate limits.}  The receipt-parsing
    pipeline uses the Gemini API, which imposes per-minute and per-day
    quotas on the free tier.  The system was designed to queue parsing
    requests gracefully rather than fail silently under burst load, and the
    economic analysis confirmed that current MES submission volumes stay
    well within the free-tier limits.
  \item \textbf{OCR accuracy.}  LLM-based extraction is not 100\% accurate
    for all receipt formats.  The design therefore preserves the raw
    receipt image and includes a manual correction step before final
    submission, allowing users to verify extracted values.
  \item \textbf{Online-only operation.}  The system has no offline mode.
    This is consistent with the rest of the MES website and acceptable for
    the administrative use case, but it means the system is unavailable
    during MES infrastructure outages.
\end{itemize}

\subsection{Assumptions}

\begin{itemize}
  \item \textbf{SSO dependency.}  All users authenticate through the
    existing MES Single Sign-On system.  The finance module does not
    manage credentials independently; it relies on the identity established
    by SSO to determine role assignments.
  \item \textbf{Timely submissions.}  The workflow assumes that club
    treasurers submit receipts within the current academic year.
    Cross-year submissions are not explicitly handled.
  \item \textbf{Schema isolation.}  The MongoDB database schema can be
    extended with new collections and fields for the finance module without
    disrupting existing MES data.  This was confirmed in the September~26
    technical meeting with the MES Infra-Tech team.
\end{itemize}

\subsection{Constraints}

\begin{itemize}
  \item \textbf{MES monorepo integration.}  The most significant
    architectural constraint was the requirement to develop within the MES
    monorepo using the slices architecture.  This precluded a standalone
    deployment and required conforming to the existing Next.js App Router
    conventions, pre-styled component library, and CI/CD pipeline.
  \item \textbf{Mandated technology stack.}  MongoDB, Next.js, and the
    MES component library were all client-mandated; no alternatives were
    evaluated for these layers.
  \item \textbf{MIT licence.}  The project inherits the MIT licence from
    the MES monorepo.  This was agreed with MES and ensures that future
    student teams and MES Infra-Tech can maintain and extend the system
    freely.
  \item \textbf{RBAC role hierarchy.}  The three-tier role model (club
    member $\rightarrow$ club executive $\rightarrow$ VP Finance/MES admin)
    was specified directly by MES stakeholder requirements and is not a
    design choice open to revision.
\end{itemize}

% ----------------------------------------------------------------
\section{Economic Considerations (LO23)}

The MES Club Payment Tracking System is an open-source contribution to a
non-profit student organisation and is not intended for commercial sale.
The relevant economic questions therefore concern sustainability, adoption,
and potential for reuse rather than pricing or profit.

\paragraph{Target users and market size.}
The McMaster Engineering Society supports approximately 6,000 undergraduate
engineering students organised into 40$+$ clubs.  Within those clubs the
system is used by roughly 40--80 club treasurers (primary submitters),
80--150 club executives (approvers), and 5--10 MES VP Finance and
administrative staff.  The total active user base in a given academic year
is therefore in the range of 200--300 people.

\paragraph{Production and operating cost.}
The system runs on the existing MES cloud infrastructure managed by the
Infra-Tech team at no additional hosting cost.  The only variable operating
cost is the Google Gemini API for receipt parsing, which at current
submission volumes falls within the free tier.  Estimated ongoing cost:
\$0--\$20 CAD per month.  If MES submission volumes grew by an order of
magnitude, a paid Gemini tier would cost on the order of \$10--\$50 per
month still negligible for a student society budget.

\paragraph{Attracting users and driving adoption.}
Because the system is embedded in the MES website, user acquisition is
automatic once the MES Infra-Tech team merges the module into the
production branch and retires the existing Google~Forms process.  The
Rev~0 demonstration and the April~2026 usability testing session were
conducted partly to build confidence with the Infra-Tech stakeholders and
accelerate that merge decision.

\paragraph{Potential for broader adoption.}
The MIT licence and the slices architecture mean that any other university
engineering society running a similar Next.js monorepo could adopt the
module with minimal effort.  Student engineering societies at other
Canadian universities face identical club-finance problems.  Attracting
broader adoption would require extracting the finance slice into a
standalone package and publishing it with setup documentation, an effort
estimated at two to four weeks of developer time beyond the capstone scope.

% ----------------------------------------------------------------
\section{Reflection on Project Management (LO24)}

\subsection{How Does Your Project Management Compare to Your Development Plan}

The team followed the Development Plan closely in most respects.  Weekly
60-minute team meetings were held consistently throughout the project.
Stakeholder check-ins with Luke (MES) occurred roughly bi-weekly as
planned.  GitHub Issues and pull requests were used for all code and
documentation changes, Discord for day-to-day coordination, and email for
formal stakeholder communication.

The two notable deviations were:
\begin{itemize}
  \item The team developed entirely within the MES monorepo
    (\texttt{MES-Website-App-Router}) rather than using a submodule
    approach.  Documentation remained in the capstone repository as
    planned, so grading and contribution tracking were not affected.
  \item The CI/CD plan was initially limited to automated LaTeX PDF
    generation.  ESLint linting, TypeScript type-checking, and the staging
    deployment pipeline were added later than planned; this gap was noted
    in the Rev~0 feedback and corrected before the final submission.
\end{itemize}

\subsection{What Went Well?}

\begin{itemize}
  \item \textbf{GitHub-based workflow.}  Feature branches, required PR
    reviews, and issue-linked commits kept every contribution traceable
    and made it easy to identify what changed and why.
  \item \textbf{Early de-risking of OCR.}  The receipt-parsing POC
    succeeded ahead of schedule and removed the largest technical
    uncertainty from the project, allowing the team to focus on the
    application layer rather than on the AI pipeline.
  \item \textbf{Consistent team meetings.}  Meeting notes were recorded as
    GitHub Issues, creating a searchable decision log for the duration of
    the project.
  \item \textbf{Structured usability testing.}  Adopting a task-based
    think-aloud methodology produced specific, actionable findings rather
    than vague impressions, and the resulting report was comprehensive
    enough to drive targeted UI improvements.
\end{itemize}

\subsection{What Went Wrong?}

\begin{itemize}
  \item \textbf{POC development lagged.}  The TA noted at Meeting~\#3
    (Oct~23) that POC implementation ``should be going, but it's not.''
    The delay was caused by underestimating the time needed to understand
    the MES monorepo and slices architecture before being able to add new
    code.
  \item \textbf{Team contribution tracking started too late.}
    Contributions were not logged consistently from day one, which led to
    an underwhelming Rev~0 productivity report (Prof.\ Smith:
    ``underwhelming productivity report this data is valuable to the TA
    and instructor'').
  \item \textbf{Miscommunication at Rev~0 demo.}  Two team members
    attended virtually due to a last-minute miscommunication, which
    affected the cohesion of the live presentation.
  \item \textbf{Instructor annotations left in deliverables.}
    \texttt{\textbackslash wss\{\}} grader comments were left in submitted
    PDFs at Rev~0, creating an unprofessional appearance.  These were
    removed in PR~\#123 before the final submission.
\end{itemize}

\subsection{What Would you Do Differently Next Time?}

\begin{itemize}
  \item \textbf{Set up contribution tracking from week one.}  A standing
    agenda item to review each member's commits, PRs, and meeting
    attendance at every weekly meeting would have prevented the Rev~0
    reporting gap.
  \item \textbf{Start implementation in parallel with early documentation.}
    Treating documentation and implementation as sequential phases caused
    the POC lag.  Future projects would benefit from a spike on the riskiest
    technical component as soon as the problem is understood, regardless of
    where documentation stands.
  \item \textbf{Pre-submission checklist.}  A shared checklist run before
    every graded submission (remove template comments, verify author names,
    confirm placeholder text is gone) would have caught the annotation and
    productivity-report issues before they reached graders.
  \item \textbf{Explicit attendance confirmation before presentations.}
    Confirming in-person vs.\ remote participation at least 24~hours
    before every graded presentation would prevent last-minute
    miscommunications.
\end{itemize}

% ----------------------------------------------------------------
\section{Ethics and Equity (LO18)}

\paragraph{Financial data privacy.}
The system handles sensitive reimbursement data belonging to student clubs.
RBAC enforcement ensures that a club member can only access their own
submissions, a club executive can access submissions within their own club,
and VP Finance/MES admin accounts have system-wide read and approval access.
No user can view another club's financial records.  All access is logged
in an immutable audit trail, providing transparency and accountability for
every approval and rejection decision.

\paragraph{Accessibility.}
The system targets the Web Content Accessibility Guidelines (WCAG) as
documented in the SRS non-functional requirements.  Pre-styled MES
components already adhere to the accessibility standards maintained by the
MES Infra-Tech team, and new components developed for the finance module
were designed to remain consistent with those standards.

\paragraph{Equitable treatment across clubs.}
All 40$+$ MES clubs access the same system with the same workflow.  No
club receives preferential treatment in the review pipeline.  The uniform
approval flow and audit trail ensure that decisions are made on the basis
of submitted documentation rather than personal relationships or manual
discretion.

\paragraph{Data minimisation.}
The system collects only the information strictly necessary for
reimbursement processing: receipt image, amount, description, and club
association.  No additional personal data is stored beyond what is already
held by the MES SSO system.

\paragraph{Universal design.}
The interface was designed to serve three distinct user types with very
different usage patterns: casual club members who may submit only one or
two receipts per year, club executives who review submissions regularly,
and VP Finance staff managing high transaction volumes.  The usability
testing confirmed that a first-time user with no prior documentation or
training could successfully complete all core workflows.

% ----------------------------------------------------------------
\section{Reflection on Capstone}

\subsection{Which Courses Were Relevant}

\begin{itemize}
  \item \textbf{SE~3XA3 (Software Engineering Practice):}  The most
    directly applicable course.  The requirements documentation process
    (SRS), hazard analysis methodology, V\&V planning and reporting
    workflow, and the concept of iterative revision based on formal
    feedback all mirror the SE~3XA3 course structure at a larger scale.
  \item \textbf{SE~3S03 / CS~3SH3 (Software Testing):}  Foundational for
    writing the VnV Plan and Report, designing test cases with explicit
    coverage criteria, and distinguishing between unit testing, system
    testing, and acceptance testing.
  \item \textbf{CS~3DB3 (Databases):}  Relevant to MongoDB schema design,
    query correctness for the financial data model, and understanding the
    implications of document-based storage (vs.\ relational) for data
    integrity in a financial context.
\end{itemize}

\subsection{Knowledge/Skills Outside of Courses}

\begin{itemize}
  \item \textbf{Next.js App Router (server components and server actions).}
    The project uses the Next.js~13$+$ App Router paradigm, including
    React Server Components, server actions, and streaming.  This
    architecture was not covered in any SE or CS course; team members
    learned it from the Next.js documentation and the existing MES codebase.
  \item \textbf{LLM API integration (Google Gemini).}  Designing a
    production-grade pipeline to extract structured financial data (line
    items, subtotals, HST) from unstructured receipt photographs using a
    generative AI API required knowledge of prompt engineering, output
    parsing, and error handling not taught in the curriculum.
  \item \textbf{Role-Based Access Control at the application layer.}
    Implementing RBAC with server-side route guards, middleware-level
    authorisation checks, and API-level permission enforcement required
    learning patterns beyond what is typically covered in SE security
    courses.
  \item \textbf{MongoDB Atlas and document-based schema design.}  Working
    within the existing MES MongoDB schema required understanding
    document modelling trade-offs (embedding vs.\ referencing) and the
    operational constraints of a shared production database.
  \item \textbf{Large-scale LaTeX documentation.}  Managing a
    multi-document LaTeX project with shared definition files, automated
    PDF generation via GitHub Actions, and cross-document
    \texttt{\textbackslash externaldocument} references was a new skill
    for most team members.
  \item \textbf{CI/CD within a large organisation's monorepo.}  Writing
    GitHub Actions workflows that integrate into an existing organisation's
    CI pipeline including linting, type checking, and staging deployment
    against a protected main branch required practical knowledge that
    goes beyond typical course assignments.
\end{itemize}

\end{document}
